\chapter{Практика 5. OFDM модуляция}

\section{Задание}
Реализовать формирование OFDM символа, включающего QPSK символы, опорные сигналы, нулевые защитные интервалы и циклический префикс.

\section{Теоретическое описание}
OFDM (Orthogonal Frequency Division Multiplexing) — метод мультиплексирования, при котором данные передаются на множестве ортогональных поднесущих. Основные элементы OFDM символа:

\begin{enumerate}
    \item \textbf{Информационные поднесущие} — передают QPSK символы
    \item \textbf{Опорные сигналы (пилоты)} — известные символы для оценки канала
    \item \textbf{Нулевые поднесущие} — защитные интервалы для борьбы с внеполосным излучением
    \item \textbf{Циклический префикс (CP)} — копия конца символа, добавляемая в начало для борьбы с межсимвольной интерференцией
\end{enumerate}

Количество опорных поднесущих:
\[
N_{RS} = \left\lfloor \frac{N_{QPSK}}{\Delta RS} \right\rfloor
\]

Количество нулевых поднесущих:
\[
N_2 = C \cdot (N_{RS} + N_{QPSK})
\]

\section{Реализация}

\begin{lstlisting}[language=Python, caption=OFDM модулятор]
class OfdmModulator:
    def __init__(self, delta_rs=6, cp_len=16, c_param=0.25, pilot_value=complex(0.707, 0.707)):
        self.delta_rs = delta_rs
        self.cp_len = cp_len
        self.c_param = c_param
        self.pilot_value = pilot_value
        
    def modulate(self, symbols):
        self.n_qpsk = len(symbols)
        self.n_pilots = int(np.floor(self.n_qpsk / self.delta_rs))
        if self.n_pilots == 0:
            self.n_pilots = 1
        self.n_zeros = int(self.c_param * (self.n_pilots + self.n_qpsk))
        self.total_carriers = self.n_pilots + self.n_qpsk + 2 * self.n_zeros
        
        spectrum = np.zeros(self.total_carriers, dtype=complex)
        
        self.pilot_indices = []
        pilot_pos = self.n_zeros
        for i in range(self.n_pilots):
            idx = pilot_pos + i * self.delta_rs
            if idx < self.total_carriers - self.n_zeros:
                self.pilot_indices.append(int(idx))
                spectrum[idx] = self.pilot_value
        
        self.data_indices = []
        data_idx = 0
        for i in range(self.total_carriers):
            if i >= self.n_zeros and i < self.total_carriers - self.n_zeros:
                if i not in self.pilot_indices:
                    if data_idx < len(symbols):
                        spectrum[i] = symbols[data_idx]
                        self.data_indices.append(i)
                        data_idx += 1
        
        time_signal = np.fft.ifft(spectrum)
        cp = time_signal[-self.cp_len:]
        ofdm_symbol = np.concatenate([cp, time_signal])
        return ofdm_symbol
\end{lstlisting}

\section{Результат}
При параметрах: $\Delta RS = 6$, $T_{CP} = 16$, $C = 0.25$, $N_{QPSK} = 285$:
\begin{itemize}
    \item $N_{RS} = 47$ опорных поднесущих
    \item $N_2 = 83$ нулевых поднесущих
    \item Общее количество поднесущих: 498
    \item Длина OFDM символа: $498 + 16 = 514$ отсчетов
\end{itemize}

\section{Вывод}
OFDM модулятор успешно реализован. Он формирует OFDM символ с учетом всех необходимых компонентов: опорные сигналы, информационные данные, нулевые поднесущие и циклический префикс.